\chapter{Introduction}
\label{chap:introduction}

Earth observation (EO) satellites now capture more imagery than ground segments can usefully downlink and process within the timeframes that matter most: a flood spreading overnight, a wildfire front shifting with the wind. Moving inference from the ground to the spacecraft removes the downlink bottleneck for exactly these time-critical applications, and a growing number of missions have begun to demonstrate that this is technically feasible \cite{barretta2026onboardaireview}. The European Space Agency's $\Phi$-sat-2 mission is among the most complete demonstrations to date: a 6U CubeSat carrying an 8-band high-resolution multispectral imager alongside the Ubotica CogniSAT-XE1 payload, built around the space-qualified Intel Movidius Myriad 2 Vision Processing Unit (VPU) \cite{myriad2}. Rather than running a single fixed application, $\Phi$-sat-2 uses the NanoSat MO Framework \cite{coelho2021phisat2} to install, update, and orchestrate multiple AI applications on board, much as a smartphone runs apps rather than a single program.

Every on-board application flown to date, however, has been a small, task-specific convolutional neural network (CNN) trained, validated, and uplinked individually for one job: cloud masking, ship detection, or a single land-cover product. Adding a new capability means repeating the entire pipeline from scratch. Geospatial Foundation Models (GeoFMs) offer a different model of development. Pre-trained on large, diverse archives of EO imagery, a GeoFM such as TerraMind \cite{jakubik2025terramind} learns a general-purpose semantic representation that downstream tasks can exploit with comparatively few labeled examples, rather than requiring a bespoke architecture and full training run per application. This property matters specifically for applications like flood and burnt-area mapping, where labeled events are rare by nature and a mission cannot simply wait to accumulate a large annotated archive before it becomes useful in an emergency.

Two obstacles stand between this promise and an operational on-board GeoFM, and they are largely independent of each other. The first is architectural: GeoFMs built on Vision Transformers (ViTs) are computationally and structurally incompatible with space-qualified accelerators of the Myriad 2's generation, which predate the hardware and software support ViT inference requires. The second is a domain gap. GeoFMs such as TerraMind are pre-trained on Sentinel-2 and similar archived missions; $\Phi$-sat-2 is a physically different sensor, and the statistical properties of its imagery differ from Sentinel-2's in ways that a model never saw during pre-training. Solving either problem alone does not solve the other: a compressed model built without addressing the domain gap simply compresses the gap along with it, and a domain-adapted ViT still cannot run on the target hardware.

This thesis addresses both obstacles together, using TerraMind and $\Phi$-sat-2 as the concrete case study, with flood segmentation as the motivating downstream application. The methodological approach combines three elements. First, a physics-based sensor degradation pipeline, adapted from the simulator developed for the ESA $\Phi$-lab OrbitalAI challenge \cite{longepe2024orbitalai}, is used to project labeled Sentinel-2 imagery into the $\Phi$-sat-2 sensor domain, producing labeled training data for a target sensor that has not itself accumulated labeled flood events. Second, a paired-contrastive domain-adaptation objective exploits a purpose-built dataset of spatially and temporally co-registered image triplets -- real $\Phi$-sat-2, real Sentinel-2, and simulated $\Phi$-sat-2, all of the same scene -- to train the model to represent the same underlying content identically regardless of which sensor produced it. Third, this domain-invariance objective is applied directly during knowledge distillation \cite{hinton2015distilling} of the TerraMind encoder into a lightweight CNN, so that the deployable student model inherits sensor-invariance as a property of how it is trained, rather than as an incidental byproduct of compression.

The data underlying this work falls into two categories. The first is the co-registered triplet dataset itself: real $\Phi$-sat-2 acquisitions, obtained via the ESA Insula Perception platform, matched to concurrent Sentinel-2 scenes and to a physics-simulated $\Phi$-sat-2 rendering of the same Sentinel-2 scene, co-registered using the LightGlue deep feature-matching algorithm \cite{Lindenberger_2023_ICCV} and filtered for spatio-temporal proximity and cloud cover. This dataset carries land-cover labels derived from ESA WorldCover 2021, providing the only currently available real-sensor ground truth. The second category is simulated: a physics-degraded, graded version of the Sen1Floods flood-segmentation benchmark, produced with the same degradation pipeline, standing in for a real labeled $\Phi$-sat-2 flood dataset that does not yet exist because the mission has not captured an annotated flood event.

\section{Problem Definition}
\label{sec:problem-definition}

The central technical problem is a covariate shift between the sensor domain a GeoFM is pre-trained on and the sensor domain it is asked to operate in. Let $X$ denote the space of satellite image patches and $Y$ the space of semantic labels (land-cover class, flood/no-flood, and so on). Writing $p_s$ and $p_t$ for the source (pre-training, e.g.\ Sentinel-2) and target (operational, e.g.\ $\Phi$-sat-2) distributions, a covariate shift is present when
\begin{equation}
    p_t(x) \neq p_s(x) \quad \text{while} \quad p_t(y \mid x) = p_s(y \mid x).
    \label{eq:covariate-shift}
\end{equation}
That is, the physical scene and its semantic content are identical across the two domains -- the same flooded field is still a flooded field -- but the pixel statistics produced by the sensor observing it differ, because of differing spectral response functions, point spread functions, signal-to-noise characteristics, and geometric band-to-band alignment. A model whose learned representation is sensitive to these sensor-specific statistics rather than to the underlying semantics will degrade on the target domain even though nothing about the task itself has changed.

Prior work on GeoFM domain gap has typically measured this shift only indirectly, through downstream task performance, and has compared source and target datasets that differ in geography, acquisition date, and scene content in addition to sensor -- confounding the sensor-specific shift this thesis is concerned with, with shifts due to the underlying data distribution itself. The co-registered triplet dataset used in this work is designed specifically to remove that confound, since the real $\Phi$-sat-2, real Sentinel-2, and simulated patches used for evaluation depict the same scene at the same time.

\section{Research Gap}
\label{sec:research-gap}

Two bodies of closely related work sharpen, rather than close, the gap this thesis addresses. Du et al.\ \cite{du2025onorbitdemonstrationgeospatialfoundation} recently demonstrated an on-orbit GeoFM (a compressed Prithvi variant) on the Kanyini CubeSat and the IMAGIN-e payload, adapting it to the target sensor by retraining task heads on real, though limited, target-domain labels while keeping the encoder frozen. Their measured domain gap, however, is confounded in the way described above: source and target are entirely different benchmark datasets, not co-registered pairs, so sensor-induced shift cannot be isolated from content and geographic shift. Separately, $\Phi$satNet, a bespoke U-Net-based GeoFM developed within the same research group specifically for on-board execution on $\Phi$-sat-2 hardware, sidesteps the ViT-compatibility problem altogether by training a compact, non-transformer architecture from scratch on a purpose-built pre-training corpus, rather than adapting an existing large-scale pre-trained model. Neither work isolates sensor-only covariate shift using co-registered data, and neither tests whether an existing GeoFM's broad, diverse pre-training transfers to a new sensor with fewer labels than a mission-specific model trained from scratch requires -- the specific comparison this thesis undertakes against $\Phi$satNet in Chapter~\ref{chap:results-discussion}.

The gap this thesis addresses is therefore threefold: (i) whether the domain gap between $\Phi$-sat-2 and Sentinel-2-pretrained TerraMind can be measured directly, in latent space, using data that isolates sensor-induced covariate shift from confounding factors; (ii) whether that gap can be closed through an explicit, paired-contrastive domain-invariance objective rather than incidental exposure to degraded training data; and (iii) whether the resulting invariance survives compression into a hardware-deployable, cross-architecture (ViT-to-CNN) student model.

\section{Research Objective and Problem Statement}
\label{sec:research-objective}

The objective of this thesis is to determine whether TerraMind can be made to produce domain-invariant latent representations under the $\Phi$-sat-2 sensor domain gap, and whether this invariance can be preserved through knowledge distillation into a lightweight CNN suitable for deployment on the $\Phi$-sat-2 on-board processor, without sacrificing the representational benefits of large-scale GeoFM pre-training. This objective is pursued through three research questions:

\begin{enumerate}
    \item Does the domain gap caused by sensor differences between $\Phi$-sat-2 and Sentinel-2 measurably affect the latent representations produced by TerraMind, and if so, where in the encoder does this effect concentrate?
    \item Can TerraMind be trained to produce more domain-invariant representations using a domain-adaptation technique that exploits co-registered, physics-simulated sensor pairs, given that naive fine-tuning on degraded data alone does not achieve this?
    \item Can the resulting domain-invariant encoder be distilled into a lightweight CNN suitable for on-board $\Phi$-sat-2 inference while preserving the learned invariance, rather than only preserving task accuracy?
\end{enumerate}

Answering these questions would show that the generalization promise of large GeoFMs -- broad semantic transfer from limited labels -- is compatible with the low-latency, resource-constrained inference that on-board emergency-response applications require, provided the domain gap is addressed deliberately rather than assumed away.
